\input {preamble}\begin {document}\fancyhead {}\fancyfoot {}
\par \noindent \textbf {\textnumero 1}\par \begin {otherlanguage*}{english}\PFUudcName {}004.729, 519.872\par \makeatletter \textbf {Palagashvili \,Abuli\unhbox \voidb@x \nobreak \,Mikhailovich}\makeatother \par \textit {Comparative Analysis of Active Queue Management Algorithms: RED, ARED, Gentle~RED}\par \textbf {Background.} Active Queue Management (AQM) algorithms of the RED family remain widely deployed in network equipment, yet a systematic three-way comparison of RED, Adaptive RED (ARED), and Gentle~RED under strictly identical conditions is lacking in the literature. \textbf {Purpose.} The aim of this study is to quantify the differences in buffer management dynamics, throughput stability, and congestion signaling among the three algorithms and to formulate practical selection guidelines. \textbf {Method.} A dumbbell-topology network model was implemented in the ns-3 discrete-event simulator (version~3.41) with ten TCP bulk-send sources creating sustained congestion at a 2~Mbit/s bottleneck link. All three algorithms were evaluated with identical queue thresholds, load profile, and topology over a 60-second simulation. An extended metric set including EWMA zone distribution and temporal drop structure was applied. \textbf {Results.} All algorithms achieved approximately 90\% link utilization (mean throughput 1.80~Mbit/s). RED provided the most predictable throughput (coefficient of variation 18.7\%) at the cost of the highest drop count (1240~packets). ARED minimized drops (1003~packets, 19\% fewer than RED) but produced the largest queue oscillations (p99 of 49~packets). Gentle~RED achieved the lowest mean queue length (6.40~packets) with the highest throughput variability (coefficient of variation 24.7\%). \textbf {Conclusions.} No single algorithm is universally optimal; the choice depends on scenario priorities: Gentle~RED for minimum delay, ARED for minimum packet loss, and RED for maximum throughput predictability. \par \keywordsname :active queue management, AQM, RED, Adaptive RED, Gentle RED, congestion control, ns-3, network simulation, bufferbloat\end {otherlanguage*}\par \par \medskip 
